\section{Separator consistency and restricted splits}
\label{sec:consistency}

An affine separator minorant is a restricted cost shift between two bags.
This section identifies the loss caused by restricting those shifts.  The
bag minima in this section are exact; an error bound for relaxed bags is
given at the end.  These results explain the consistency part of a
certificate.  The construction and verification of finite certificates
also require the local models and compatibility checks of
\cref{def:model-certificate}.

\begin{definition}[Separator splits]
\label{def:consistency-split}
Let $N$ be the number of bags.  Identify a nonroot bag $t$ with its edge
to its parent $\pi(t)$.  A split is a family of bounded functions
$\chi_t:X_{S_t}\to\R$, $t\ne r$.  Its shifted bag functions and bound are
\begin{equation}
\label{eq:consistency-split}
 a_t^\chi(z)=a_t(z)+\sum_{u\in\mathrm{ch}(t)}\chi_u(z_{S_u})
                  -\chi_t(z_{S_t}),
 \qquad
 \rho(\chi)=\sum_t\inf_{z\in X_{V_t}}a_t^\chi(z),
\end{equation}
where the last term in $a_r^\chi$ is omitted.
For linear spaces $\Phi_t$ of bounded separator functions containing the
constants, put $\Phi=\prod_{t\ne r}\Phi_t$ and
\[
 \rho(\Phi)=\sup_{\chi\in\Phi}\rho(\chi),\qquad
 \operatorname{gap}(\Phi)=f^*-\rho(\Phi).
\]
Write $\operatorname{gap}(\chi)=f^*-\rho(\chi)$ for a fixed split.
Constants in a split telescope, and $\sum_ta_t^\chi(x_{V_t})=F(x)$.
Thus every split gives a valid bound and every gap is nonnegative.
\end{definition}

For an edge $t$, let $A_t=\sub(t)$ and let $B_t$ be the other bags.
Define the value functions of the two sides by
\[
 U_t(s)=\min_{x_{S_t}=s}\sum_{v\in A_t}a_v(x_{V_v}),\qquad
 V_t(s)=\min_{x_{S_t}=s}\sum_{v\in B_t}a_v(x_{V_v}).
\]
The child-side function $U_t$ is the conditional value $\phi_t$ of
\cref{eq:model-value}.  The minimizations use the variables appearing on
the respective side.
Running intersection makes their private variables disjoint.  Hence
\begin{align}
\label{eq:consistency-band}
 L_t&=f^*-V_t,\qquad w_t=U_t-L_t\ge0,\qquad
 f^*=\min_s(U_t(s)+V_t(s)),\nonumber\\
 \mathcal B_t&=\{\psi:X_{S_t}\to\R\text{ bounded}:L_t\le\psi\le U_t\}.
\end{align}
The functions $U_t,V_t$ are continuous: each is a minimum of a continuous
function over a fixed compact product of private-variable intervals.
The margin $w_t(s)$ also equals
$\min\{F(x)-f^*:x_{S_t}=s\}$, and therefore vanishes at every optimal
separator value.  Denote by $\mathcal E$ the set of exact bounded splits,
that is, the splits with $\rho(\psi)=f^*$.

The dynamic-programming split $\psi_t=U_t$ belongs to $\mathcal E$.
Indeed its nonroot shifted bag is
$a_t+\sum_{u\in\mathrm{ch}(t)}U_u-U_t$, which is nonnegative and has
infimum zero by the defining recursion for $U_t$.  The root shifted bag
has infimum $f^*$.  The same argument works for bounded discontinuous bag
functions if minima are replaced by infima.  This observation will allow
us to absorb bounded splits when considering either side of an edge.

\subsection{Duality and bag envelopes}

The dual interpretation of restricted splits is agreement of neighboring
bag measures on a chosen family of test functions.  We give the compact
form needed here. The proof uses compact minimax \citep{Sion1958} and
allows compact convex sets of normalized functionals.

\begin{theorem}[Restricted split duality]
\label{thm:consistency-duality}
Suppose $\Phi_t\subset C(X_{S_t})$ for every edge.  For each bag let
$\mathcal M_t$ be a nonempty weak-star compact convex subset of
$C(X_{V_t})^*$ whose members satisfy $\ell(1)=1$.  Then
\begin{align}
\label{eq:consistency-duality}
 &\sup_{\chi\in\Phi}\sum_t\min_{\ell\in\mathcal M_t}\ell(a_t^\chi)
 \nonumber\\
 &\quad=
 \min\left\{\sum_t\ell_t(a_t):
 \ell_t\in\mathcal M_t,\quad
 \ell_t(h)=\ell_{\pi(t)}(h)
 \text{ for all }h\in\Phi_t,\ t\ne r\right\}.
\end{align}
Separator functions in the constraints are pulled back to the two bags.
When the constraint set is nonempty, its minimum is attained.  When it is
empty, both sides have value $+\infty$.
For $\mathcal M_t$ equal to the probability measures on $X_{V_t}$, the
left side is $\rho(\Phi)$.
\end{theorem}

\begin{proof}
On the compact convex product $\mathcal M=\prod_t\mathcal M_t$, consider
\[
 Q(\ell,\chi)=\sum_t\ell_t(a_t)
 +\sum_{t\ne r}\bigl(\ell_{\pi(t)}(\chi_t)-\ell_t(\chi_t)\bigr).
\]
For a fixed split this is continuous affine in $\ell$.  For fixed $\ell$
it is continuous linear in the separator functions equipped with their
uniform norms.  The compact minimax theorem gives
$\sup_\chi\min_\ell Q=\min_\ell\sup_\chi Q$.
The supremum of the mismatch terms is zero if all the stated constraints
hold, and is $+\infty$ otherwise, because the spaces are linear.
The consistent families form a closed subset of $\mathcal M$, proving
attainment when they exist.  For probability measures, minimizing the
expectation of a continuous bag function is the same as minimizing the
function itself.
\end{proof}

Containment of all point evaluations in each $\mathcal M_t$ is sufficient
for the functional bound to be a valid relaxation of the bag minima.
Normalization alone does not ensure this; the general identity in
\cref{thm:consistency-duality} is distinct from validity of a chosen bag
relaxation.

If all separator tests are allowed, neighboring separator marginals are
equal.  Such bag measures have a joint extension on the tree.  To see
this, sample the root bag from its measure.  For each child, sample its
remaining variables from a conditional law given the separator.
Conditional laws exist on compact metric spaces.  Running intersection
ensures that those remaining variables have not yet been assigned and
that the separator already has the required law.  Induction gives every
prescribed bag marginal.  The expected objective of this joint law is at
least $f^*$, confirming exactness of unrestricted splitting.

\begin{proposition}[Closed bag envelopes]
\label{prop:consistency-envelope}
For a bounded function $f$ on a compact box, the infimum of $f$ equals
the minimum of its closed convex envelope.
Let $h_t^\chi$ be the closed convex envelope of $a_t^\chi$ on $X_{V_t}$.
For any spaces of bounded split functions containing constants,
\begin{equation}
\label{eq:consistency-envelope}
 \sup_{\chi\in\Phi}\min_{x\in X}\sum_t h_t^\chi(x_{V_t})
 =\rho(\Phi+\mathrm{Aff}),
\end{equation}
where $\mathrm{Aff}$ consists of affine separator functions.
In particular,
$\rho(\mathrm{Aff})=\min_{x\in X}\sum_t\operatorname{vex}_{X_{V_t}}a_t(x_{V_t})$.
\end{proposition}

\begin{proof}
Fix a split and write $f_t=a_t^\chi$.  The set
\[
 K_t=\cl\conv\{(z,f_t(z)):z\in X_{V_t}\}
\]
is compact.  Its lower fiber boundary
$h_t(z)=\min\{v:(z,v)\in K_t\}$ is lower semicontinuous, convex, and
below $f_t$.  Every closed convex minorant of $f_t$ has a closed convex
epigraph containing $K_t$, so it is below $h_t$.  Thus $h_t=h_t^\chi$.
The infimum of a linear functional over a set equals its infimum over
the closed convex hull.  Consequently, for every affine $b$,
\[
 \min_{(z,v)\in K_t}(v+b(z))=\inf_z(f_t(z)+b(z)).
\]
The case $b=0$ proves the individual bag statement.

Minimize $\sum_tv_t$ over $(z_t,v_t)\in\prod_tK_t$ with equal neighboring
separator copies.  Running intersection identifies these consistent
copies with one global $x$, and minimization in the fibers gives
$\min_x\sum_th_t(x_{V_t})$.
Apply compact minimax to this compact product and unrestricted affine
separator multipliers.  The supremum over the multipliers enforces
equality of the separator copies.  The preceding linear-functional
identity then gives
\[
 \min_x\sum_th_t(x_{V_t})
 =\sup_{b\in\mathrm{Aff}}\sum_t\inf a_t^{\chi+b}.
\]
Taking the supremum over $\chi$ proves \cref{eq:consistency-envelope}.
The argument uses boundedness, not continuity, of the chosen split.
For the last assertion take $\Phi$ to consist of constants.
\end{proof}

\begin{proposition}[Attainment and affine exactness]
\label{prop:consistency-attainment}
If every $\Phi_t$ is a finite-dimensional subspace of
$C(X_{S_t})$, then $\rho(\Phi)$ is attained.
Consequently, $\operatorname{gap}(\mathrm{Aff})=0$ if and only if an
affine split has a global point whose bag projections minimize all its
shifted bags.  For an exact split, every global minimizer has this property.
\end{proposition}

\begin{proof}
Quotient out constant edge functions, which leave the split objective
unchanged.  If the quotient is zero-dimensional, attainment is immediate.
Otherwise, for a direction $d$ in the remaining finite-dimensional
space, put $D_t=\sum_{u\in\mathrm{ch}(t)}d_u-d_t$, omitting $d_t$ at the
root.  Since $\sum_tD_t(x_{V_t})=0$,
$\sum_t\min D_t\le0$.
If equality holds, all the nonnegative differences
$D_t(x_{V_t})-\min D_t$ vanish for every global $x$.
Every bag point extends to a point of the box, so every $D_t$ is constant.
Starting at the leaves then shows that every $d_t$ is constant.
Thus the sum is strictly negative for every nonzero quotient direction.
It is continuous on the quotient unit sphere, so it is at most
$-\kappa$ there for some $\kappa>0$.  For a direction of norm $R$,
\[
 \rho(d)\le\sum_t\sup a_t+\sum_t\inf D_t
 \le\sum_t\sup a_t-\kappa R.
\]
The split objective is continuous, because changing the split uniformly
changes every bag infimum by at most the sum of the incident uniform
changes.  Its upper level sets are therefore compact on the quotient,
which proves attainment.

If one point minimizes every shifted bag, the sum of these minima is
its objective value.  Validity forces that value to be $f^*$.
Conversely, for an exact split and a global minimizer $x^*$, the
nonnegative differences $a_t^\chi(x^*_{V_t})-\inf a_t^\chi$ sum to zero,
so every difference vanishes.
\end{proof}

\subsection{One separator and the need for joint control}

For bounded functions, distances below use the uniform norm.  For a
function class $\Phi$ and a set $\mathcal B$ write
$\dist_\infty(\Phi,\mathcal B)=\inf_{\chi\in\Phi,\psi\in\mathcal B}
\norm{\chi-\psi}_\infty$, and write
$\dist_\infty(f,\Phi)=\inf_{\chi\in\Phi}\norm{f-\chi}_\infty$.

\begin{theorem}[The one-separator band identity]
\label{thm:consistency-band}
Consider two sides joined by one separator, with value functions $U,V$
and $L=f^*-V$.  Let $\mathcal B=\{L\le\psi\le U\}$ and let $\Phi$ be a
linear space of bounded separator functions containing constants.  Then
\begin{equation}
\label{eq:consistency-band-distance}
 \operatorname{gap}(\Phi)
 =\inf_{\chi\in\Phi}\left\{\sup(\chi-U)+\sup(L-\chi)\right\}
 =2\dist_\infty(\Phi,\mathcal B).
\end{equation}
For continuous $\Phi,U,L$, the band in the distance can be restricted to
continuous functions.
If the separator has a finite single-valued partition $\mathcal P$ and
$\Phi$ allows independent pieces $\chi_D\in\Phi_D$ on each cell, where
every $\Phi_D$ is a linear space containing constants, then
\begin{equation}
\label{eq:consistency-worst-cell}
 \operatorname{gap}(\Phi)=\max_{D\in\mathcal P}g_D,\qquad
 g_D=\inf_{\chi_D\in\Phi_D}
 \left\{\sup_D(\chi_D-U)+\sup_D(L-\chi_D)\right\}.
\end{equation}
Individual cell brackets $g_D$ may be negative.
\end{theorem}

\begin{proof}
The two side infima for a split $\chi$ are
$\inf(U-\chi)=-\sup(\chi-U)$ and
$\inf(V+\chi)=f^*-\sup(L-\chi)$, proving the bracket expression.
Put $a=\sup(\chi-U)$ and $b=\sup(L-\chi)$.
We have $a+b\ge-\inf(U-L)=0$.
Shift $\chi$ by $(b-a)/2$, so both suprema become $(a+b)/2$.
Clipping the shifted function into the band gives
$\psi=\min\{\max\{\chi,L\},U\}$ and
$\norm{\chi-\psi}_\infty\le(a+b)/2$.
This proves the lower inequality for the distance identity; the clip is
continuous when its three input functions are continuous.

Conversely, for $\psi\in\mathcal B$,
\[
 \sup(\chi-U)+\sup(L-\chi)
 \le\sup(\chi-\psi)+\sup(\psi-\chi)=\osc(\chi-\psi).
\]
For any bounded function $q$,
$\inf_{c\in\R}\norm{q-c}_\infty=\osc(q)/2$.
Constants belong to the class, so optimizing gives the other inequality.

A global bracket is at least the bracket of its restriction to each
cell.  For the reverse inequality, choose a piece on each cell within
$\delta$ of $g_D$, and shift that piece by a constant so its two
suprema are equal.  The resulting global bracket is at most
$\max_Dg_D+\delta$.  Let $\delta$ decrease to zero.
\end{proof}

Thus a restricted split need not approximate a particular value function:
it needs to approximate some function between $L$ and $U$.  At separator
values where the margin is positive, this band provides room for error.
For a zero-width band, the gap is twice the usual uniform approximation
error.  With piecewise constants, the cell bracket is
$\sup_D L-\inf_D U$.

\begin{theorem}[Tree sandwich]
\label{thm:consistency-tree}
For the tree and split classes of \cref{def:consistency-split} with
$N\ge2$,
\begin{equation}
\label{eq:consistency-tree}
 2\max_{t\ne r}\dist_\infty(\Phi_t,\mathcal B_t)
 \le\operatorname{gap}(\Phi)
 \le2\inf_{\psi\in\mathcal E}\sum_{t\ne r}\dist_\infty(\psi_t,\Phi_t)
 \le2\sum_{t\ne r}\dist_\infty(U_t,\Phi_t).
\end{equation}
The infimum in the middle is over one joint exact split.
\end{theorem}

\begin{proof}
For $\psi\in\mathcal E$ write $m_t=\inf a_t^\psi$, so
$\sum_tm_t=f^*$.  If $q_t=\chi_t-\psi_t$ has infimum $\alpha_t$ and
supremum $\beta_t$, then
\[
 \inf a_t^\chi\ge m_t+\sum_{u\in\mathrm{ch}(t)}\alpha_u-\beta_t,
\]
with no last term at the root.  Summing gives
$\operatorname{gap}(\chi)\le\sum_{t\ne r}\osc(q_t)$.
Constants in each class allow independent centering, so optimizing the
oscillations gives the middle bound.  The last bound uses the exact split
$U$.

For the lower bound fix one edge and enlarge every other class to all
bounded functions.  For a fixed function on the chosen edge, exact
dynamic programming on each side reduces its optimized bound to
$\inf(U_t-\chi_t)+\inf(V_t+\chi_t)$.
This remains valid for the bounded modified bag functions, by the
infimum-based recursion given above.  Enlarging classes can only increase
the bound.  Apply \cref{thm:consistency-band} to this two-side problem
and then maximize over the chosen edge.
\end{proof}

\begin{example}[Original bands do not control a joint affine split]
\label{ex:consistency-obstruction}
On $[-1,1]^2$ take a three-bag path, rooted at an endpoint, with functions
\[
 a_r(s_1)=2\beta s_1^2,\qquad
 a_m(s_1,s_2)=\kappa(s_1-s_2)^2-\beta s_2^2,\qquad
 a_l(s_2)=2\beta s_2^2,
\]
where $\beta>0$ and $\kappa\ge2\beta$.
Here $f^*=0$ and both original bands contain zero, but
$\operatorname{gap}(\mathrm{Aff})=\beta$.
Indeed the first band's endpoints are
\[
 U_1(s_1)=\frac{\kappa\beta}{\kappa+\beta}s_1^2,
 \qquad L_1(s_1)=-2\beta s_1^2,
\]
and the second has $U_2(s_2)=2\beta s_2^2$ and
\[
 L_2(s_2)=-\left(\frac{2\beta\kappa}{2\beta+\kappa}-\beta\right)s_2^2\le0.
\]
These follow by minimizing quadratics whose unconstrained minimizers lie
in the box.  After removing telescoping constants, an affine split is
$\chi_i=\lambda_i s_i$.  Evaluating the middle shifted bag at
$(s_1,s_2)=(1,1)$ and $(-1,-1)$ bounds its infimum by
$-\beta-\abs{\lambda_2-\lambda_1}$.  The other two infima are at most zero
by evaluation at zero.  Hence $\rho(\mathrm{Aff})\le-\beta$.
The zero split attains $-\beta$, proving the claim.
In particular, refinement driven only by the affine brackets of the
original edge bands can stop without reaching the desired joint gap.
\end{example}

\subsection{Graded exact splits and discounted bag errors}

A particular joint exact split retains a controlled part of every edge
band.  This gives a margin-sensitive bound in every separator dimension.

\begin{theorem}[Graded exact splits]
\label{thm:consistency-graded}
Assume $n_T=N-1\ge1$.  For each nonroot bag $t$, let $b_t$ be the number of
edges strictly inside $\sub(t)$, and put
\begin{equation}
\label{eq:consistency-graded}
 \tau=\frac1{2n_T},\qquad
 \theta_t=(2b_t+1)\tau,\qquad \psi_t=U_t-\theta_tw_t.
\end{equation}
Then $\psi$ is exact.  For all bounded perturbations $q_t$,
\begin{equation}
\label{eq:consistency-discount}
 \operatorname{gap}(\psi+q)
 \le\sum_{t\ne r}\left\{\sup(q_t-\tau w_t)+\sup(-q_t-\tau w_t)\right\}.
\end{equation}
Define the narrower band
$\mathcal C_t=\{\sigma:\abs{\sigma-\psi_t}\le\tau w_t\}$.
It lies in $\mathcal B_t$, and
$\operatorname{gap}(\Phi)\le2\sum_{t\ne r}\dist_\infty(\Phi_t,\mathcal C_t)$.
For independent cell classes on partitions $\mathcal P_t$, with every
$\Phi_{t,D}$ a linear space containing constants, this gives
\begin{equation}
\label{eq:consistency-cell-sandwich}
 \operatorname{gap}(\Phi)\le\sum_{t\ne r}\max_{D\in\mathcal P_t}g_{t,D},
 \quad
 g_{t,D}=\inf_{\ell\in\Phi_{t,D}}
 \left\{\sup_D(\ell-\psi_t-\tau w_t)
              +\sup_D(\psi_t-\tau w_t-\ell)\right\}.
\end{equation}
\end{theorem}

\begin{proof}
Define the dynamic-programming bag margins
\[
 g_t=a_t+\sum_{u\in\mathrm{ch}(t)}U_u-U_t\quad(t\ne r),\qquad
 g_r=a_r+\sum_{u\in\mathrm{ch}(r)}U_u-f^*,
\]
and the bag-projection margin
\begin{equation}
\label{eq:consistency-bag-margin}
 W_t(z)=\min\{F(x)-f^*:x_{V_t}=z\}.
\end{equation}
Fixing a bag separates its child subtrees and its parent side.  Therefore
$W_t=g_t+w_t$ for $t\ne r$, and $W_r=g_r$.
Furthermore $W_t(z)$ dominates every incident separator margin evaluated
at its corresponding projection, because fixing the bag fixes each such
separator.

Write $\chi_t=U_t-e_t$, where $e_t=\theta_tw_t-q_t$.
The gap is the sum of bag deficits
\[
 D_t=\sup_z\left\{\sum_{u\in\mathrm{ch}(t)}e_u(z_{S_u})
                         -e_t(z_{S_t})-g_t(z)\right\},
\]
omitting $e_t$ at the root.  For a nonroot bag the expression inside this
supremum is
\begin{align*}
 &\sum_{u\in\mathrm{ch}(t)}(-q_u-\tau w_u)+(q_t-\tau w_t)\\
 &\quad+\sum_{u\in\mathrm{ch}(t)}(\theta_u+\tau)w_u
                  +(1-\theta_t+\tau)w_t-W_t.
\end{align*}
All coefficients of margins in the second line are nonnegative.
The identity $b_t=\sum_{u\in\mathrm{ch}(t)}(b_u+1)$ implies
$\theta_t=\sum_u(\theta_u+\tau)+\tau$, so their sum is one.
Since each margin is at most $W_t$, that line is nonpositive.
At the root, $\sum_u(\theta_u+\tau)=1$, giving the same conclusion with
no own-separator term.  Bound the remaining terms by separate suprema
and sum over the bags.  Each edge appears twice, proving
\cref{eq:consistency-discount}.

Taking $q=0$ makes the right side zero because every margin vanishes at
an optimal separator value.  Thus $\psi$ is exact.
We have $\tau\le\theta_t\le1-\tau$, so $\mathcal C_t\subset\mathcal B_t$.
For $\sigma_t\in\mathcal C_t$ and $\chi_t\in\Phi_t$, the edge bracket
in \cref{eq:consistency-discount} is at most
$\osc(\chi_t-\sigma_t)$.  Centering by constants and optimizing gives
the distance bound.  On a finite cell partition, choose near-minimizing
pieces and balance their two suprema independently, as in
\cref{thm:consistency-band}; the edge bracket is then at most its largest
cell bracket.  This proves \cref{eq:consistency-cell-sandwich}.
\end{proof}

For one edge, $\theta=\tau=1/2$ and $\mathcal C$ is the full band, so
\cref{eq:consistency-discount} is the exact one-edge bracket.
For many edges, the grading provides room on both sides of every split.
The scalar construction in \cref{sec:scalar-covering} uses this room to
bound separator cell counts by coverings of projected near-optimal sets.
It assumes exact bag minima and scalar separators; it does not give a
general bound on the number of certificate leaves.

\begin{proposition}[Discounted bag errors]
\label{prop:consistency-bag-errors}
Let $n_T=N-1\ge1$, and let $b_t$ be the number of edges strictly inside
$\sub(t)$ as in \cref{thm:consistency-graded}.
Let $\widetilde a_t=a_t-\mathrm{err}_t$, where
$\mathrm{err}_t:X_{V_t}\to\R$ is bounded and nonnegative, and let
$\widetilde\rho$ be its split bound with exact bag infima.  Using the
original margins $w_t,W_t$, put
\[
 \tau'=\frac1{3n_T+1},\qquad
 \theta'_t=(3b_t+2)\tau',\qquad \psi'_t=U_t-\theta'_tw_t.
\]
Then for all bounded perturbations $q$,
\begin{equation}
\label{eq:consistency-bag-errors}
 f^*-\widetilde\rho(\psi'+q)
 \le\sum_{t\ne r}\left\{\sup(q_t-\tau'w_t)+\sup(-q_t-\tau'w_t)\right\}
       +\sum_t\sup_{z\in X_{V_t}}\left(\mathrm{err}_t(z)-\tau'W_t(z)\right).
\end{equation}
\end{proposition}

\begin{proof}
The bag deficit in the preceding proof gains $\mathrm{err}_t$.
Reserve the term $\mathrm{err}_t-\tau'W_t$ and repeat the decomposition
with $\tau',\theta'_t$.  For a nonroot bag, the remaining margin
coefficients are $\theta'_u+\tau'$ for its children and
$1-\theta'_t+\tau'$ for its own separator.  They are nonnegative and
their sum is $1-\tau'$, because
$\theta'_t=\sum_u(\theta'_u+\tau')+2\tau'$.
They multiply margins bounded by $W_t$, and are accompanied by
$-(1-\tau')W_t$.  Their total is therefore nonpositive.
At the root $\sum_u(\theta'_u+\tau')=3n_T\tau'=1-\tau'$, which gives the
same inequality.  Sum the separate suprema as before.
\end{proof}

\begin{remark}[Computational scope]
\label{rem:consistency-scope}
The value functions and projection margins are global quantities.
The preceding approximation bounds do not supply inexpensive oracles
for them.  Independent separator pieces also require evaluating a bag on
the common refinement of its incident partitions.  The product of the
incident cell counts is an upper bound on the number of such regions;
overlapping separator coordinates can reduce it.  Consequently a small
number of separator coefficients alone is not a bound on computational
work.  The following sections construct and verify finite affine
certificates directly, with their stated local oracles and error budgets.
\end{remark}
